\documentclass[11pt,a4paper]{ctexart}

\usepackage[margin=2.2cm]{geometry}
\usepackage{amsmath}
\usepackage{booktabs}
\usepackage{longtable}
\usepackage{xcolor}
\usepackage[hidelinks]{hyperref}

\setlength{\parindent}{2em}
\setlength{\parskip}{0.35em}
\linespread{1.12}
\sloppy

\newcommand{\code}[1]{\texttt{\detokenize{#1}}}
\newcommand{\bps}{\mathrm{bps}}

\title{glftmm\_lot 策略说明}
\author{基于当前仓库代码与配置}
\date{2026-06-16}

\begin{document}
\maketitle
\tableofcontents
\newpage

\section{阅读入口}

这份文档描述当前仓库里的 \code{SimpleMakerStrategy}，以及它在
\code{config/config.json} 下形成的策略族。主要代码入口是：

\begin{itemize}
  \item \code{run_strategy.py}：读取配置、展开参数组合、恢复每日状态、写出 parquet 和 summary。
  \item \code{sim/strategy.py}：核心做市策略、profit grid、止损、超时、multi-lot 状态机。
  \item \code{sim/loader.py}：把 sampler 缓存和原始成交流合成时间序列事件。
  \item \code{core/manager.py}、\code{core/orderbook.py}、\code{core/position.py}：本地订单簿、撮合、仓位和 PnL。
  \item \code{sampler/*.py}：采样盘口、交易强度、波动率、instructor 信号。
\end{itemize}

一句话概括：这是一个事件驱动的被动做市回测。策略在固定频率的盘口采样点更新双边报价，用原始成交流模拟自己的 maker/taker 成交；持仓盈利时改为 close-only 的分层止盈网格。multi-lot 版本把仓位组织成从内到外的 lot 栈，内层只允许平仓，只有最外层可开仓 lot 才按原单 lot 规则自由开平仓。

\section{当前配置下的策略族}

当前 \code{config/config.json} 的交易对象是 \code{BTCUSDT}，日期为
\code{2025-01-02} 到 \code{2026-06-01}。策略读取原始数据根目录
\code{input_path} 和 \code{input_backup_path}，读取 sampler 缓存根目录
\code{output_path=/data/users/kang/backtest/glftmm_var/cached}，策略结果写到
\code{/data/users/kang/backtest/glftmm_lot/result}。

本文描述策略运行时以 \code{config/config.json} 为准。\code{sampler/config.json}
当前自己的 \code{output_path} 是 \code{/data/users/kang/backtest/glftmm_lot/cached}；
如果两份配置不同，策略实际读取的是 \code{config/config.json} 里的缓存路径。

当前会展开成 16 个策略任务：

\begin{itemize}
  \item \code{mode} 有 2 个值：\code{1} 和 \code{0}。
  \item \code{profit_grid} 有 4 组。
  \item \code{max_position_usdt} 和 \code{stoploss} 成对，有 2 组 lot 结构。
  \item 其他参数当前都是单值。
\end{itemize}

\begin{longtable}{p{0.31\linewidth}p{0.60\linewidth}}
\toprule
配置项 & 当前含义 \\
\midrule
\endhead
\code{mode=[1,0]} & 撮合难度。0 是 inclusive，1 是 strict。mode=1 更难成交。\\
\code{latency=2} & ticker 事件生成的 maker 报价在 2 毫秒后才进入本地 maker book。\\
\code{price_precision=1} & 价格 tick 为 0.1。\\
\code{qty_precision=3} & 数量 step 为 0.001。\\
\code{taker_fee=0.00025} & taker 成交扣 2.5 bps。\\
\code{maker_fee=-0.00003} & maker 成交是负手续费，即返佣 0.3 bps。\\
\code{freq=1000} & 每 1000 ms 读取一次采样盘口和特征。\\
\code{name_intensity} / \code{lookback_intensity} & 当前为 \code{k} / \code{300}，用 300 个采样点的破价强度均值作为基础报价距离。\\
\code{name_volatility} / \code{lookback_volatility} & 当前为 \code{sigma} / \code{300}，读取 sigma 波动率，但当前 \code{adj_spread_volatility=0}，不改变报价。\\
\code{name_instructor} / \code{lookback_instructor} & 当前为 \code{bbo_imbalance} / \code{0}，读取当前采样盘口深度不平衡，但当前 \code{adj_spread_instructor=0}，不改变报价。\\
\code{open_passive_only=true} & 如果将来启用 instructor，只允许它把普通开仓报价变得更被动，不允许追价开仓。当前因 instructor 缩放为 0，实际无影响。\\
\code{order_amt=1000} & 每次基础单边下单名义金额约 1000 USDT，再按数量精度和最小下单约束截断。\\
\code{min_order_notional=100} & 开仓部分必须满足最小名义金额。\\
\code{inventory_skew=[0,1]} & 最大 skew tick 为 0，所以当前不开库存偏移。\\
\code{open_curve_underwater} & 当前为 \code{[0,1,0]}，order 为 0，所以当前不开水下补仓数量衰减。\\
\code{new_open_lot_crit=0.01} & 下一层 lot 只有在上一层满仓且价格相对上一层成本价逆向偏离 1\% 后才允许开仓。\\
\code{max_holding_time=-1} & 不启用持仓超时强平。\\
\code{profit_grid} & 四组止盈网格，分别测试不同盈利带和数量分布。\\
\bottomrule
\end{longtable}

当前 lot 结构如下：

\begin{itemize}
  \item 五层 lot：每层 \code{max_position_usdt=100000}，每层 \code{stoploss=30000}，总仓位上限 500000 USDT。
\end{itemize}

\section{数据流和时间顺序}

策略运行前需要 sampler 缓存。当前策略本身不负责生成 sampler 缓存；它只读取已经存在的 parquet。如果缺少某一天的采样盘口、强度、波动率或 instructor 文件，策略会直接报错。

\subsection{采样盘口}

\code{sampler/resample.py} 把原始 \code{BOOKTICKER} 转成固定网格。每个采样时间戳使用该时刻及之前最后一条已知盘口：

\[
\text{sampled\_ticker}(t)=\text{last bookticker update at or before }t.
\]

它不会使用 \([t,t+\text{freq})\) 内未来才出现的盘口。输出列包括：
\code{timestamp}、\code{best_bid_price}、\code{best_ask_price}、
\code{best_bid_qty}、\code{best_ask_qty}。

\subsection{交易强度}

当前策略使用 \code{name_intensity=k}。它的基础量是每个桶内成交价突破当时 BBO 的最大距离：

\[
\text{break\_dist}
=\max(\max(\text{trade\_price}-\text{best\_ask},0),
      \max(\text{best\_bid}-\text{trade\_price},0)).
\]

代码先对 \code{bucket_break_dist_max} 做 \code{shift(1)}，所以采样点 \(t\) 使用的是上一桶 \([t-\text{freq},t)\) 已经发生的破价信息。指标 \code{k} 是非零破价距离的滚动均值；\code{k_decay} 是指数衰减；\code{k_median} 是非零破价距离滚动中位数。

\subsection{波动率}

\code{sampler/volatility.py} 先构造每个采样桶的 mid、open、high、low、close，然后支持多种波动率：
\code{sigma}、\code{atr}、\code{rv}、\code{bv}、\code{parkinson}、\code{gk}。当前策略读取 \code{sigma}，但由于 \code{adj_spread_volatility=0}，它不会参与报价距离。

\subsection{Instructor}

支持两类 instructor：

\begin{itemize}
  \item \code{trade_imbalance}：按成交方向聚合上一窗口的买卖量，先 \code{shift(1)}，再计算
  \[
  \frac{\text{bid\_qty}-\text{ask\_qty}}
       {\text{bid\_qty}+\text{ask\_qty}+\epsilon}.
  \]
  这个信号严格不使用同桶成交。
  \item \code{bbo_imbalance}：当前采样盘口深度不平衡：
  \[
  \frac{\text{best\_bid\_qty}-\text{best\_ask\_qty}}
       {\text{best\_bid\_qty}+\text{best\_ask\_qty}+\epsilon}.
  \]
  它只支持 \code{lookback=0}。
\end{itemize}

当前配置读取 \code{bbo_imbalance}，但 \code{adj_spread_instructor=0}，所以它不改变实际报价。

\subsection{事件合并}

\code{BinanceEventLoader} 读取三类数据：采样 ticker/alpha 行、原始 trade 行、可选的强度/波动率/instructor 列。它输出两种事件：

\begin{itemize}
  \item \code{trade}：\code{timestamp, is_buyer_maker, price, volume}。
  \item \code{ticker}：\code{timestamp, best_bid, best_ask, instructor, intensity, volatility}。
\end{itemize}

如果 trade 和 ticker 时间戳相同，trade 先输出。策略处理每个事件前会先激活已经到期的 pending maker quote。因此，若一张挂单的 \code{active_ts} 正好等于某笔 trade 的时间戳，它可以先进入本地订单簿，再参与这笔 trade 的撮合；该时间戳的 ticker 更新会在 trade 之后处理。

\section{单 lot 策略的报价公式}

每个 ticker 事件先检查 BBO 是否有效：bid、ask 必须有限且为正，并且 bid < ask。中间价为：

\[
mid_t=\frac{best\_bid_t+best\_ask_t}{2}.
\]

基础报价距离为：

\[
d_t = intensity_t \cdot adj\_spread\_intensity
    + volatility_t \cdot adj\_spread\_volatility.
\]

当前配置下，\code{adj_spread_intensity=1} 且 \code{adj_spread_volatility=0}，所以：

\[
d_t=intensity_t.
\]

Instructor 的方向位移为：

\[
shift_t = instructor_t \cdot adj\_spread\_instructor \cdot mid_t.
\]

如果 \code{open_passive_only=true}，普通开仓报价会做 clipping：

\[
shift^{ask}_t=\max(0,shift_t),\qquad
shift^{bid}_t=\min(0,shift_t).
\]

这意味着 instructor 只能让 ask 往上、bid 往下，即让开仓更被动；不会让 ask 往下或 bid 往上追价。当前配置 \code{adj_spread_instructor=0}，所以 \(shift_t=0\)。

库存偏移为：

\[
skew_t =
\begin{cases}
-\lfloor |u_t|^{p}\cdot max\_skew\_ticks \rfloor\cdot tick, & u_t>0,\\
\phantom{-}\lfloor |u_t|^{p}\cdot max\_skew\_ticks \rfloor\cdot tick, & u_t<0,\\
0, & u_t=0,
\end{cases}
\]

其中 \(u_t=cost\_notional\_usdt/max\_position\_usdt\)。多头库存会把双边报价往下移，鼓励卖出、抑制继续买入；空头库存会把双边报价往上移。当前 \code{max_skew_ticks=0}，所以不开库存偏移。

最终普通报价为：

\[
ask\_price = ceil\_tick(best\_ask+d_t+shift^{ask}_t+skew_t),
\]

\[
bid\_price = floor\_tick(best\_bid-d_t+shift^{bid}_t+skew_t).
\]

基础单边数量为：

\[
q = floor\_step(order\_amt / mid_t).
\]

如果 \(q\) 小于 \code{min_order_qty} 或 \code{min_order_notional/mid_t}，则该侧数量变成 0。当前 \code{order_amt=1000}，\code{min_order_notional=100}，所以一般情况下基础订单不会被最小名义金额挡住。

\section{库存上限和普通开仓}

单 lot 的仓位上限使用 \code{cost_notional_usdt} 判定，不使用 mark notional：

\[
cost\_notional\_usdt = qty \cdot cost.
\]

规则是：

\begin{itemize}
  \item 如果 \(cost\_notional\_usdt \ge max\_position\_usdt\)，停止继续挂 bid 开多。
  \item 如果 \(cost\_notional\_usdt \le -max\_position\_usdt\)，停止继续挂 ask 开空。
  \item 如果 \code{max_position_usdt<=0}，双边数量都为 0。
\end{itemize}

注意这里是方向性 cost notional。多头达到上限时不再买，但仍可卖；空头达到上限时不再卖，但仍可买。

\section{挂单延迟、订单簿和撮合}

Ticker 事件生成的 maker levels 不会立即永久替换订单簿，而是先进入 pending 队列：

\[
active\_ts = timestamp + latency.
\]

每个新事件开始时，策略会把所有已到期的 pending quote 弹出，只保留最后一条到期报价，清空旧 maker book，然后放入这条报价。这样做的效果是：如果市场事件很稀疏，中间过期的旧报价不会逐条执行，只有最新到期报价生效。

\subsection{mode=0 和 mode=1}

\code{core/orderbook.py} 里 \code{mode} 决定成交阈值：

\begin{itemize}
  \item ask 挂单：mode=0 时 \(ask\_price \le trade\_price\) 成交；mode=1 时 \(ask\_price < trade\_price\) 才成交。
  \item bid 挂单：mode=0 时 \(bid\_price \ge trade\_price\) 成交；mode=1 时 \(bid\_price > trade\_price\) 才成交。
\end{itemize}

所以 mode=1 比 mode=0 更难成交，尤其是在 trade price 正好等于挂单价时，mode=1 不会填。

\subsection{成交方向}

原始 trade 的 \code{is_buyer_maker=true} 表示买方是 maker，主动方是卖方，也就是市场卖单打到了被动 bid。策略注释中直接写了这一点。对应关系是：

\begin{itemize}
  \item \code{is_buyer_maker=true}：公共卖流打 bid，本策略的 bid maker 可能成交，仓位增加。
  \item \code{is_buyer_maker=false}：公共买流打 ask，本策略的 ask maker 可能成交，仓位减少或开空。
\end{itemize}

每笔 trade 先尝试撮合 maker book。若成交量还有剩余，再撮合 taker book。maker 成交用 \code{maker_fee}，taker 成交用 \code{taker_fee}。

\section{profit grid 止盈状态机}

当前策略开启 \code{profit_grid}，所以持仓后的行为不是继续双边普通报价，而是根据持仓盈亏切换：

\begin{itemize}
  \item 空仓：正常双边开仓。
  \item 多头且 \(mid \le cost\)：仍处于水下或刚回本，只挂 bid 继续补多，不挂 ask 止盈。
  \item 多头且 \(mid > cost\)：只挂 ask close-only 止盈网格，不再挂 bid 开仓。
  \item 空头且 \(mid \ge cost\)：仍处于水下或刚回本，只挂 ask 继续补空，不挂 bid 止盈。
  \item 空头且 \(mid < cost\)：只挂 bid close-only 止盈网格，不再挂 ask 开仓。
\end{itemize}

也就是说，profit grid 把单 lot 的行为分成两个区域：

\begin{itemize}
  \item 水下区域：允许沿持仓方向继续补仓，但受仓位上限和 underwater 曲线影响。
  \item 盈利区域：只做 close-only 分层止盈。
\end{itemize}

\subsection{网格价格}

\code{profit_grid=[lower_bps, upper_bps, num_grid, qty_kind, qty_param]}。多头止盈价格在：

\[
cost\cdot(1+lower/10000)
\quad \text{到} \quad
cost\cdot(1+upper/10000)
\]

之间均匀取 \code{num_grid} 档，并向上取 tick。空头止盈价格在：

\[
cost\cdot(1-lower/10000)
\quad \text{到} \quad
cost\cdot(1-upper/10000)
\]

之间均匀取档，并向下取 tick。

如果多头某个止盈价低于当前 \code{best_ask}，会被抬到 \code{best_ask}；如果空头某个止盈价高于当前 \code{best_bid}，会被压到 \code{best_bid}。这避免 close-only 网格挂在明显不合理的位置。

\subsection{网格数量分布}

可选数量分布有三类：

\begin{itemize}
  \item \code{equal}：每档同权。
  \item \code{power}：权重为 \((num\_grid-index)^{param}\)，越靠近第一档权重越大。
  \item \code{exponential}：权重为 \(\exp(-param\cdot index)\)，param 越大，越集中在第一档。
\end{itemize}

代码先按权重分配总平仓数量，再按 step 向下取整。若某档数量小于最小合法下单量，则后续档不再分配。最后剩余数量会加到第一档。因此第一档通常承担最多数量，特别是在 exponential 分布下。

当前配置的四组网格是：

\begin{itemize}
  \item \code{[5,105,10,"exponential",0.2]}：5 到 105 bps，10 档，数量缓慢向第一档集中。
  \item \code{[5,105,10,"exponential",0.5]}：同样盈利带，数量更集中在第一档。
  \item \code{[2,52,10,"exponential",0.2]}：更近的 2 到 52 bps，10 档。
  \item \code{[2,22,5,"exponential"]}：2 到 22 bps，5 档，exponential 默认参数为 1。
\end{itemize}

\subsection{超过上界后的释放}

如果 mid 已经超过 profit grid 上界，策略不再维持整条网格，而是生成一档更容易成交的 close-only 单：

\begin{itemize}
  \item 多头：价格为 \(\min(ceil(best\_bid+tick), best\_ask)\)。
  \item 空头：价格为 \(\max(floor(best\_ask-tick), best\_bid)\)。
\end{itemize}

这不是市价单，但比原网格更接近盘口，用于在价格已经穿过目标盈利区后尽快释放持仓。

\subsection{重复网格保护}

close-only 网格有一个 \code{profit_grid_key}，由 ask 价格序列、bid 价格序列和持仓成本组成。如果新 ticker 计算出的 close-only 网格和上一条完全相同，策略不会清空并重挂 maker book，只记录状态。这减少了重复改单。

\section{水下补仓数量曲线}

当 profit grid 开启且持仓在水下时，补仓数量会通过 \code{open_curve_underwater} 乘一个系数：

\[
util=\min(1,\max(0,gross\_cost\_notional\_usdt/max\_position\_usdt)),
\]

\[
multiplier=min\_scale+(max\_scale-min\_scale)\cdot(1-util)^{order}.
\]

如果 \code{order<=0}，代码直接返回 1。当前配置为 \code{[0,1,0]}，所以当前不会随仓位利用率降低补仓数量。

如果将来把 order 调成正数，仓位越接近该 lot 的上限，补仓数量越小；当 util 接近 1 时，multiplier 接近 \code{min_scale}。

\section{止损和 reach-and-release}

单 lot 的止损按未实现浮亏触发：

\[
floating\_unrealized = qty\cdot(mid-cost),
\]

\[
floating\_loss=\max(0,-floating\_unrealized).
\]

如果 \code{floating_loss >= stoploss}，策略进入 \code{reach_and_release_active}。进入该状态后：

\begin{itemize}
  \item 清空 pending maker quote。
  \item 清空 maker book。
  \item 清空旧 taker book。
  \item 按当前仓位方向放一张 close-only taker 单。
\end{itemize}

多头止损时，策略在 ask taker book 放卖出单，价格用最新 \code{best_bid}；空头止损时，策略在 bid taker book 放买入单，价格用最新 \code{best_ask}。它不是在 ticker 事件当场直接改仓，而是等后续 trade 事件来撮合 taker book。成交后用 \code{taker_fee} 计费。

当仓位完全归零后，\code{reach_and_release_active} 清除，超时跟踪也清除。

\section{最大持仓时间}

\code{max_holding_time} 控制达到单 lot 仓位上限后的超时释放：

\begin{itemize}
  \item \code{max_holding_time<0}：关闭该逻辑。当前配置为 -1。
  \item \code{max_holding_time=0}：只要已经到达上限并记录了开始时间，下一个 ticker 就触发释放。
  \item \code{max_holding_time>0}：达到上限后，如果当前时间戳减开始时间大于该值，则触发释放。
\end{itemize}

只要仓位被减小，超时开始时间会重置。这里的“达到上限”同样用 \code{gross_cost_notional_usdt >= max_position_usdt} 判定。

\section{multi-lot 分层逻辑}

\code{glftmm_lot} 和普通单 lot 策略的主要区别在这里。配置中的每个 \code{max_position_usdt} 元素都会生成一个 \code{_SingleLotMakerStrategy}，并配对一个同位置的 \code{stoploss}。这些子策略共享同一条市场事件流，但仓位、成本、订单簿、止盈网格、止损状态分别保存。

当前 multi-lot 不是“多个 lot 各自独立交易”。它被组织成一个从内到外的逻辑栈：

\begin{itemize}
  \item index 小的是内层 lot，index 大的是外层 lot。
  \item \code{active_lot_count} 表示当前开放的前缀长度；\code{active_lot_indices} 是 \code{0..active_lot_count-1}。
  \item \code{frozen_lot_indices} 是尚未开放的外层备用 lot。
\end{itemize}

\subsection{初始状态}

多 lot 初始化时：

\begin{itemize}
  \item \code{active_lot_count=1}。
  \item \code{active_lot_indices=[0]}，只有最内层 lot 开始工作。
  \item 其他 lot 在 \code{frozen_lot_indices} 中，作为外层备用 lot。
\end{itemize}

每个事件开头，策略会先激活各 lot 已经到期的 pending maker quote。但 trade 事件不会再让所有 lot 各自独立吃同一笔成交；它只在 active 前缀内，从内到外依次撮合，并把上一层已经成交的数量从这笔 trade 中扣掉。

\subsection{什么时候启用下一层 lot}

每个 ticker 事件前，策略先同步 lot 栈。如果当前 active 前缀的最外层 lot 已经有仓位，并且满足：

\begin{itemize}
  \item 该 lot 已经达到自己的仓位上限；
  \item 判断标准是 \code{gross_cost_notional_usdt >= max_position_usdt}；
  \item 当前价格相对该 lot 成本价达到 \code{new_open_lot_crit} 的逆向偏离；
  \item 多头满仓时要求 \code{mid <= cost * (1 - new_open_lot_crit)}，空头满仓时要求 \code{mid >= cost * (1 + new_open_lot_crit)}；
  \item 这个价格门槛只控制下一层是否能开仓，不改变当前 lot 的 stoploss 和 close-only 流程；
\end{itemize}

则 \code{active_lot_count} 加 1，下一层外侧备用 lot 被开放。这个开放发生在 ticker 事件上；如果某一笔 trade 刚好让当前外层 lot 达到满仓，下一层会在后续 ticker 同步时检查满仓和价格偏离是否同时满足。

因此，多 lot 不是同时开所有层，而是只有上一层达到任意方向满仓，并且行情继续朝该层不利方向移动足够远后才启用下一层。这个机制把加仓风险拆成一层一层向外展开的仓位单元，并降低单边行情里快速打满所有 lot 的速度。

\subsection{谁可以开仓，谁只能平仓}

ticker 事件上，multi-lot 会先找当前可开仓 lot：

\begin{itemize}
  \item 如果 active 前缀里没有任何仓位，则第 0 层可以按原单 lot 规则双边开仓。
  \item 如果已经有非空 lot，通常只有最外层非空 lot 可以按原单 lot 规则自由开平仓。
  \item 如果最外层非空 lot 已经满仓、价格偏离条件仍然满足，并且外侧存在一个刚开放但仍为空的备用 lot，则这个空备用 lot 是当前可开仓 lot。
  \item 其他更内层的非空 lot 不允许开仓，只能保留现有规则下的减仓侧报价。例如多头只保留 ask 平仓侧，空头只保留 bid 平仓侧。
\end{itemize}

这意味着内层 lot 在水下时不会继续补仓；它要么等待现有规则产生可平仓报价，要么进入 stoploss / reach-and-release 的 close-only 流程。外层 lot 承担新的开仓和补仓动作。

\subsection{trade 成交优先级}

一笔 trade 到来时，active 前缀按逻辑 index 从小到大处理，也就是从最内层开始：

\[
remaining\_qty_0 = trade\_qty.
\]

第 \(i\) 层只最多撮合 \(\text{remaining\_qty}_i\)。如果该层成交了 \(filled_i\)，下一层只看到：

\[
remaining\_qty_{i+1}=remaining\_qty_i-filled_i.
\]

因此，同一笔市场成交会优先满足最内层的平仓单，再轮到次内层，以此类推。只有内层没有成交或成交后还有剩余数量，外层开仓单才可能吃到这笔 trade。

\subsection{内层平空后的 index 轮换}

同步 lot 栈时，如果 active 前缀里某个内层 lot 已经空仓，而它外侧仍有非空 lot，则这个空 lot 会被移动到 active 前缀最外层的位置。它外侧的非空 lot 逻辑 index 依次前移。

例如 active 前缀原来是：

\[
[lot_0=\text{空},\ lot_1=\text{非空},\ lot_2=\text{非空}]
\]

同步后变成：

\[
[lot_1=\text{非空},\ lot_2=\text{非空},\ lot_0=\text{空备用}]
\]

这样可以保证非空仓位总是从内到外连续排列，内层平仓优先级也随之更新。所有 frozen 且空仓的 lot 会清空 pending maker quote、maker book、taker book 和重复网格 key。

\subsection{聚合输出}

每日输出不是逐 lot 分开写 parquet，而是把所有 lot 聚合成一条策略曲线：

\begin{itemize}
  \item \code{position}：各 lot 数量求和。
  \item \code{mark_notional_usdt}：各 lot mark notional 求和。
  \item \code{cost_notional_usdt}：各 lot cost notional 求和。
  \item \code{gross_cost_notional_usdt}：各 lot gross cost notional 求和。
  \item \code{realized_pnl} 和 \code{unrealized_pnl}：各 lot 求和。
  \item \code{traded_volume}：各 lot 成交额求和。
\end{itemize}

状态 JSON 里会保留 \code{lots}，包含每一层 lot 的独立状态。这里的 \code{lots} 顺序是当前逻辑顺序，而不一定是最初配置顺序。状态也会保存 \code{active_lot_count}、\code{active_lot_indices}、\code{frozen_lot_indices}、\code{max_position_lots}、\code{stoploss_lots} 和 \code{new_open_lot_crit}，用于断点续跑和输出复盘。

\section{仓位和 PnL}

\code{Position.execute} 使用 signed quantity：

\begin{itemize}
  \item 正数表示买入，仓位增加。
  \item 负数表示卖出，仓位减少或转为空。
\end{itemize}

同方向加仓会更新平均成本。反方向成交会先平已有仓位并产生 realized PnL；如果反向成交超过原仓位，剩余部分以成交价作为新方向的成本。手续费直接从累计 realized PnL 中扣除。由于当前 maker fee 为负数，maker 成交会增加 realized PnL。

未实现 PnL 为：

\[
unrealized\_pnl = mark\_notional\_usdt - cost\_notional\_usdt.
\]

其中：

\[
mark\_notional\_usdt=qty\cdot mid,\qquad
cost\_notional\_usdt=qty\cdot cost.
\]

\section{每日运行、恢复和提前停止}

\code{run_strategy.py} 以“每日”为增量单位运行。每个参数组合的输出目录形如：

\[
\code{<output_dir>/<symbol>/<sim_dir>/strat__all}
\]

每一天会写两个文件：

\begin{itemize}
  \item \code{YYYY-MM-DD.parquet}：该日内每个 ticker 事件后的聚合曲线。
  \item \code{_state/YYYY-MM-DD.json}：该日结束后的可恢复状态。
\end{itemize}

如果 \code{overwrite=false}，已经同时存在 parquet 和 state 的日期会跳过；策略会从第一个缺失日期继续跑，并用上一日 state 恢复。

策略有三类提前停止规则，都是在 state 中记录：

\begin{itemize}
  \item \code{stop_due_to_dayn_pnl}：运行天数大于 60 且 total PnL 小于 0。
  \item \code{stop_due_to_low_annualized_return}：运行天数大于 60 且 realized PnL 年化收益率低于 0.2。
  \item \code{stop_due_to_max_drawdown}：相对历史最高 total PnL 的回撤除以总 \code{max_position_usdt} 小于等于 -0.2。
\end{itemize}

这里的年化收益率为：

\[
\frac{realized\_pnl}{max\_position\_usdt}\cdot \frac{365}{day\_n}.
\]

这里的 \code{max_position_usdt} 对多 lot 来说是所有 lot 上限之和。

\section{输出字段怎么读}

每日 parquet 字段如下：

\begin{longtable}{p{0.30\linewidth}p{0.62\linewidth}}
\toprule
字段 & 含义 \\
\midrule
\endhead
\code{timestamp} & ticker 事件时间戳，毫秒。\\
\code{price} & 当前 mid price。\\
\code{position} & 聚合仓位数量，正数为多，负数为空。\\
\code{mark_notional_usdt} & 仓位按当前 mid 标记的方向性名义金额。\\
\code{cost_notional_usdt} & 仓位按平均成本计算的方向性名义金额。\\
\code{gross_cost_notional_usdt} & \code{abs(cost_notional_usdt)}。\\
\code{realized_pnl} & 已实现 PnL，包含手续费。\\
\code{unrealized_pnl} & 未实现 PnL。\\
\code{traded_volume} & 策略自身累计成交额，所有 lot 求和。\\
\bottomrule
\end{longtable}

\code{strategy_summary.csv} 汇总每个参数组合的最终指标，包括 total PnL、realized/unrealized PnL、最终仓位、成交额、按成交额折算的 PnL bps 等。需要注意：如果某个组合提前停止，summary 反映的是停止日状态；如果历史 summary 和当前配置不匹配，应优先看输出目录下的 \code{_state}。

\section{当前配置下真正生效的机制}

当前配置里，实际改变交易行为的主要是：

\begin{itemize}
  \item \code{mode=0/1}：影响成交难度。
  \item \code{latency=2}：影响挂单进入 book 的时刻。
  \item \code{intensity=k, lookback=300, adj_spread_intensity=1}：决定普通报价离 BBO 的距离。
  \item \code{order_amt=1000} 和最小下单约束：决定普通开仓数量。
  \item \code{max_position_usdt} 和 \code{stoploss} 的多 lot 向量：决定分层仓位与每层止损。
  \item \code{new_open_lot_crit=0.01}：决定下一层 lot 需要相对上一层成本价逆向偏离 1\% 才能开仓。
  \item \code{profit_grid}：决定盈利区的 close-only 分层止盈。
  \item 每日状态恢复和三类提前停止规则：决定长区间批量回测是否继续。
\end{itemize}

当前配置里，被读取但不会改变报价的机制是：

\begin{itemize}
  \item \code{bbo_imbalance} instructor：因为 \code{adj_spread_instructor=0}。
  \item \code{sigma} 波动率：因为 \code{adj_spread_volatility=0}。
  \item \code{inventory_skew}：因为最大 tick 为 0。
  \item \code{open_curve_underwater}：因为 order 为 0。
  \item \code{max_holding_time}：因为值为 -1。
\end{itemize}

\section{几个容易误读的边界}

\begin{itemize}
  \item \textbf{profit grid 不是止损。} 它只在持仓盈利时挂 close-only 网格。亏损时只有当前可开仓 lot 仍可能补仓；内层非空 lot 只允许平仓。
  \item \textbf{止损不是立即改仓。} 止损触发后放 taker close-only 单，仍需要后续 trade 事件撮合。
  \item \textbf{多 lot 不是简单提高单层 max position，也不是多个独立策略。} 每一层有独立成本、独立订单簿、独立止损；但开仓权限、成交数量、下一层价格门槛和 index 轮换由外层 lot 栈统一调度。
  \item \textbf{仓位上限用成本名义金额。} 它不随 mark price 波动直接变化，除非成交改变平均成本或数量。
  \item \textbf{mode=1 更保守。} 等价价格不成交，只有严格穿过才成交。
  \item \textbf{当前 instructor 不生效。} 文档里解释了 instructor 公式，但当前参数缩放为 0。
  \item \textbf{latency 后只保留最后一条到期报价。} 高频 ticker 下旧 pending quote 会被新 quote 覆盖，不是每条过期报价都进入 book。
\end{itemize}

\section{读结果时建议看的顺序}

阅读某个结果目录时，建议按下面顺序：

\begin{enumerate}
  \item 先看路径里的 \code{sim__...} 参数 token，确认它是否属于当前配置族。
  \item 看最后一个 \code{_state/*.json}，确认 \code{day_n}、\code{total_pnl}、三个 stop 标记、\code{max_pnl_ever}。
  \item 如果是 multi-lot，展开 state 里的 \code{lots}，看当前逻辑顺序下每层 lot 的仓位、成本、止损状态、\code{active_lot_count} 和 active/frozen 列表。
  \item 再看每日 parquet 曲线，重点看 \code{position}、\code{cost_notional_usdt}、\code{gross_cost_notional_usdt}、\code{realized_pnl}、\code{unrealized_pnl}。
  \item 最后再用 \code{strategy_summary.csv} 排名。summary 适合批量比较，但排查机制时不如 state 直接。
\end{enumerate}

\end{document}
